\documentclass[10pt,a4paper]{amsart}
\usepackage[margin=19mm]{geometry}
\usepackage{amsmath,amssymb,mathtools}
\usepackage[scr=rsfs,cal=euler]{mathalfa}
\usepackage{xfrac}
\usepackage[unicode,hidelinks,bookmarks=false]{hyperref}
\newcommand{\ie}{i.e.\ }
\newcommand{\cf}{cf.\ }
\newcommand{\resp}{respectively\ }
\newcommand{\Subsection}{\subsection}
\providecommand{\nobreakdash}{\nobreak\hbox{-}}
\setlength{\emergencystretch}{2em}
\multlinegap0pt
\usepackage{amssymb}
\usepackage{algorithm}\usepackage{algpseudocode}
\usepackage[shortlabels]{enumitem}
\usepackage{tikz}
\newtheorem{theorem}{Theorem}
\title{\textbf{Proofs of the Conjectures on $SOME(n)$ and $DSOME(n)$ Functions Related to Integer Partitions}}
\author{Gaurab Bardhan, Nipen Saikia}
\date{Source: arXiv:2608.08524v1 (CC BY 4.0)}
\begin{document}
\maketitle

\noindent\emph{Source and licence.} \textbf{Proofs of the Conjectures on $SOME(n)$ and $DSOME(n)$ Functions Related to Integer Partitions}. Version arXiv:2608.08524v1 (CC BY 4.0), distributed by arXiv under the Creative Commons Attribution 4.0 International licence, http://creativecommons.org/licenses/by/4.0/, as recorded in the arXiv licence metadata for this exact version and shown on the abstract page https://arxiv.org/abs/2608.08524v1. Full licence notice: \texttt{licenses/arxiv-2608.08524-CC-BY-4.0.txt}.\par\smallskip

\noindent\textit{Editorial scope.} Counted unit: the article's theorem DSM8THM (source0.tex lines 1265--1277). The article's complete original proof of it is delivered with it (source0.tex lines 1279--1545, one \texttt{proof} environment of 267 lines). No mathematics is written, repaired or re-derived: the statement and the proof are the article's own lines, in the article's own notation, and no retained line is substituted or re-flowed.  Retained uncounted as the source-local context the proof needs: lines 136--138; lines 200--203; lines 287--293; lines 343--346; lines 775--782.  Nothing was omitted from any retained span, so the package carries no omission record.  No retained line contains a citation, so the wrapper carries no bibliography and no bibliography entry.  No retained line contains a figure environment, an image-inclusion command or drawing code, so no image is delivered and the package declares no asset dependency.  Editorial formatting only, in addition: an AMS \texttt{amsart} preamble that restates the article's own declarations and macros used by the retained lines, namely the environments \texttt{theorem} on separate counters, together with the standard packages those lines need.  The retained environments print in this document as Theorem 1 in order of appearance, whereas the article prints the same items with the numbering of its own full document.  The complete native source of the article as distributed in the arXiv e-print for this exact version is bundled unchanged as \texttt{sources/207169/source0.tex}.
\begin{equation}\label{DSM8M25}
\sum_{n=0}^{\infty}p_{d}(n)q^n=\dfrac{(q^2;q^2)_\infty}{(q;q)_\infty}=\dfrac{f_2}{f_1}.   
\end{equation}

\begin{equation}\label{DSE3}
\phi(-q)=\dfrac{f_1^2}{f_2}
=1+2X(q),
\end{equation}

    \begin{equation}\label{DSM8M17}
\dfrac{x(1+x^2)}{(1+x)^4}
=
\sum_{j=1}^{\infty}
(-1)^{j-1}
\dfrac{j(j^2+2)}{3}x^j.
    \end{equation}

   \begin{equation}\label{DSM8M20}
       \sum_{j\mid n}(-1)^{j-1}j^s=
\sigma_s(n)-2^{s+1}\sigma_s(n/2),
   \end{equation}

\begin{equation}\label{DSE2}
\sum_{n=0}^{\infty}DSOME(n)q^n
=
\dfrac{f_1}{8}
\left(
\phi(-q)^{-1}-\phi(-q)^3
\right).
\end{equation}

\begin{theorem}\label{DSM8THM}
 For any positive integer $n$, we have
\begin{equation*}
DSOME(n)
\equiv
3\sum_{r=0}^{n-1}
p_{\mathrm d}(r)
\left(
\sigma_3(n-r)+2\sigma(n-r)
\right)
\pmod 8.
\end{equation*}
\end{theorem}

\begin{proof}
Using \eqref{DSE2} and \eqref{DSE3}, we obtain
\begin{align}
\sum_{n=0}^{\infty}DSOME(n)q^n
=
\dfrac{f_1}{8}
\left(
\phi(-q)^{-1}-\phi(-q)^3
\right)
=
\dfrac{f_2}{f_1}\,
\dfrac{1}{8}
\left(
1-
\left(
\dfrac{f_1^4}{f_2^2}
\right)^2
\right).
\label{DSM8M3}
\end{align}
For every positive integer $m$, let
\begin{equation}\label{DSM8M4}
a_m=\dfrac{q^m}{(1+q^m)^2}.
\end{equation}
Then
$$\dfrac{f_1^4}{f_2^2}=
\dfrac{\displaystyle\prod_{m=1}^{\infty}(1-q^m)^4}
     {\displaystyle\prod_{m=1}^{\infty}(1-q^{2m})^2}=
\prod_{m=1}^{\infty}
\left(
1-\dfrac{4q^m}{(1+q^m)^2}
\right)=
\prod_{m=1}^{\infty}(1-4a_m).$$
Therefore,
\begin{equation}\label{DSM8M6}
\left(
\dfrac{f_1^4}{f_2^2}
\right)^2=
\prod_{m=1}^{\infty}(1-4a_m)^2=
\prod_{m=1}^{\infty}
\left(
1-8a_m+16a_m^2
\right).
\end{equation}
Also, for any  integer $M\geq2$, 
\begin{equation}\label{DSM8M7}
\prod_{m=1}^{M}
\left(
1-8a_m+16a_m^2
\right)
=
1+
\sum_{m=1}^{M}
\left(
-8a_m+16a_m^2
\right)\quad+
\sum_{t=2}^{M}
\sum_{1\leq m_1<\cdots<m_t\leq M}
\prod_{i=1}^{t}
\left(
-8a_{m_i}+16a_{m_i}^2
\right),
\end{equation} and for any integer 
 $t\geq2$, 
\begin{equation}\label{DSM8M8}
\prod_{i=1}^{t}
\left(
-8a_{m_i}+16a_{m_i}^2
\right)=
\prod_{i=1}^{t}
8\left(
-a_{m_i}+2a_{m_i}^2
\right)
=
8^t
\prod_{i=1}^{t}
\left(
-a_{m_i}+2a_{m_i}^2
\right)
\equiv0\pmod{64}.
\end{equation}
So employing \eqref{DSM8M8} in \eqref{DSM8M7}, we obtain
\begin{equation}\label{DSM8M9}
\prod_{m=1}^{M}
\left(
1-8a_m+16a_m^2
\right)
\equiv
1+
\sum_{m=1}^{M}
\left(
-8a_m+16a_m^2
\right)
=
1-8\sum_{m=1}^{M}a_m
+16\sum_{m=1}^{M}a_m^2
\pmod{64},
\end{equation}
Now, taking limit as
$M\rightarrow\infty$ in \eqref{DSM8M9} and employing
\eqref{DSM8M6}, we obtain
\begin{equation*}
1-
\left(
\dfrac{f_1^4}{f_2^2}
\right)^2
\equiv
8\sum_{m=1}^{\infty}a_m
-16\sum_{m=1}^{\infty}a_m^2
\pmod{64},
\end{equation*}
which implies
\begin{equation}\label{DSM8M11}
\dfrac{1}{8}
\left(
1-
\left(
\dfrac{f_1^4}{f_2^2}
\right)^2
\right)
\equiv
\sum_{m=1}^{\infty}a_m
-
2\sum_{m=1}^{\infty}a_m^2
\pmod8.
\end{equation}
From  \eqref{DSM8M4}, we note that 
\begin{align}
a_m-2a_m^2
=
\dfrac{q^m}{(1+q^m)^2}
-
\dfrac{2q^{2m}}{(1+q^m)^4}
=
\dfrac{q^m(1+q^{2m})}{(1+q^m)^4}.
\label{DSM8M12}
\end{align}
Employing  \eqref{DSM8M11} and
\eqref{DSM8M12} in \eqref{DSM8M3}, we obtain
\begin{align}
\sum_{n=0}^{\infty}DSOME(n)q^n
&\equiv
\dfrac{f_2}{f_1}
\left(
\sum_{m=1}^{\infty}a_m
-
2\sum_{m=1}^{\infty}a_m^2
\right)
\equiv
\dfrac{f_2}{f_1}
\sum_{m=1}^{\infty}
\dfrac{q^m(1+q^{2m})}{(1+q^m)^4}
\pmod8.
\label{DSM8M13}
\end{align}
Substituting $x=q^m$ in \eqref{DSM8M17}, we obtain
\begin{align}
\sum_{m=1}^{\infty}
\dfrac{q^m(1+q^{2m})}{(1+q^m)^4}
&=
\dfrac{1}{3}
\sum_{m=1}^{\infty}
\sum_{j=1}^{\infty}
(-1)^{j-1}j(j^2+2)q^{mj}
\notag\\
&=
\dfrac{1}{3}
\sum_{N=1}^{\infty}
\left(
\sum_{j\mid N}
(-1)^{j-1}j(j^2+2)
\right)q^N
\notag\\
&\label{DSM8M19}=
\dfrac{1}{3}
\sum_{N=1}^{\infty}
\left(
\sum_{j\mid N}(-1)^{j-1}j^3
+
2\sum_{j\mid N}(-1)^{j-1}j
\right)q^N.
\end{align}
Taking $s=3$ and $s=1$ in \eqref{DSM8M20}, we obtain
\begin{align}
\sum_{j\mid N}(-1)^{j-1}j^3
=
\sigma_3(N)-16\sigma_3(N/2)\quad 
\text{and}\quad 
\sum_{j\mid N}(-1)^{j-1}j
=
\sigma(N)-4\sigma(N/2).
\label{DSM8M21}
\end{align}
Employing \eqref{DSM8M21} in \eqref{DSM8M19}, we obtain
\begin{align}
3
\left(
\sum_{m=1}^{\infty}
\dfrac{q^m(1+q^{2m})}{(1+q^m)^4}
\right)
&=
\sum_{N=1}^{\infty}\left(\sigma_3(N)-16\sigma_3(N/2)+
2\left(
\sigma(N)-4\sigma(N/2)
\right)\right)q^N
\notag\\
&\equiv
\sum_{N=1}^{\infty}\left(\sigma_3(N)+2\sigma(N)\right)q^N
\pmod8,\notag
\end{align}
which gives
\begin{equation}\label{DSM8M24}
\sum_{m=1}^{\infty}
\dfrac{q^m(1+q^{2m})}{(1+q^m)^4}
\equiv
3\sum_{N=1}^{\infty}
\left(
\sigma_3(N)+2\sigma(N)
\right)q^N
\pmod8.
\end{equation}
Employing  \eqref{DSM8M25} and \eqref{DSM8M24} in \eqref{DSM8M13}, we obtain
\begin{align}
\sum_{n=0}^{\infty}DSOME(n)q^n
&\equiv
3
\left(
\sum_{r=0}^{\infty}p_{\mathrm d}(r)q^r
\right)
\left(
\sum_{N=1}^{\infty}\sigma_3(N)q^N
\right)+
6
\left(
\sum_{r=0}^{\infty}p_{\mathrm d}(r)q^r
\right)
\left(
\sum_{N=1}^{\infty}\sigma(N)q^N
\right)
\pmod8
\notag\\
&\equiv
3\sum_{n=1}^{\infty}
\left(
\sum_{r=0}^{n-1}
p_{\mathrm d}(r)\sigma_3(n-r)
\right)q^n+
6\sum_{n=1}^{\infty}
\left(
\sum_{r=0}^{n-1}
p_{\mathrm d}(r)\sigma(n-r)
\right)q^n
\pmod8
\notag\\
&\label{DSM8M26}\equiv
3\sum_{n=1}^{\infty}
\left(
\sum_{r=0}^{n-1}
p_{\mathrm d}(r)
\left(
\sigma_3(n-r)+2\sigma(n-r)
\right)
\right)q^n
\pmod8.
\end{align}
Comparing the coefficients of $q^n$ in \eqref{DSM8M26}, we complete the proof.
\end{proof}

\end{document}
